% Luc Destombe --- licence CC BY-NC-SA 4.0
% https://creativecommons.org/licenses/by-nc-sa/4.0/deed.fr
% Fichier autonome : compiler avec pdflatex, deux fois (signets, nombre de pages).
\documentclass[11pt,a4paper]{article}
\makeatletter
\newif\ifeleve
\elevefalse
\elevetrue

\newif\iflycee@palatino
\lycee@palatinotrue

\RequirePackage[utf8]{inputenc}
\RequirePackage[T1]{fontenc}
\RequirePackage[french]{babel}
\RequirePackage{etoolbox}

\newcommand{\ShorthandsOff}{\@ifundefined{shorthandoff}{}{\shorthandoff{;:!?}}}
\newcommand{\ShorthandsOn}{\@ifundefined{shorthandon}{}{\shorthandon{;:!?}}}
\ShorthandsOff
\AtBeginDocument{\ShorthandsOn}
\AtBeginEnvironment{tikzpicture}{\ShorthandsOff}

\RequirePackage{geometry}
\RequirePackage{fancyhdr}
\RequirePackage{lastpage}
\RequirePackage{multicol}
\RequirePackage{array}
\RequirePackage{enumitem}
\RequirePackage{xcolor}
\RequirePackage{needspace}

\RequirePackage{amsmath,amssymb}
\RequirePackage{mathpazo}
\DeclareMathSizes{10.95}{10.95}{8}{5.5}
\begingroup\catcode`\_=\active \catcode`\^=\active
\gdef_#1{\sb{\scriptscriptstyle #1}}
\gdef^#1{\sp{\scriptscriptstyle\lycee@moinsepais #1}}
\endgroup
\newcommand\lycee@moins{\mathbin{%
  \setbox\z@\hbox{$\m@th\scriptscriptstyle\lycee@moinsorig$}%
  \dimen@=.55\wd\z@ \advance\dimen@-.5pt
  \hbox to.75\wd\z@{\hss$\m@th\scriptscriptstyle\vcenter{\hbox{%
    \tikz\draw[line width=.5pt,line cap=round](0,0)--(\the\dimen@,0);}}$\hss}}}
\begingroup\lccode`\~=`\- \lowercase{\endgroup\let~}\lycee@moins
\newcommand\lycee@moinsepais{\mathcode`\-="8000 }
\AtBeginDocument{\mathchardef\lycee@moinsorig=\mathcode`\-\relax}
\AtBeginDocument{\catcode`\_=12 \mathcode`\_="8000
  \catcode`\^=12 \mathcode`\^="8000 }
\newcommand\lycee@exposants{%
  \fontdimen13\textfont2=.48\fontdimen6\textfont2
  \fontdimen14\textfont2=.48\fontdimen6\textfont2
  \fontdimen15\textfont2=.48\fontdimen6\textfont2
  \fontdimen13\scriptfont2=.48\fontdimen6\scriptfont2
  \fontdimen14\scriptfont2=.48\fontdimen6\scriptfont2
  \fontdimen15\scriptfont2=.48\fontdimen6\scriptfont2}
\every@math@size\expandafter{\the\every@math@size\lycee@exposants}
\newlength{\retraitcalculs}
\setlength{\retraitcalculs}{2em}

\RequirePackage{tikz}
\usetikzlibrary{arrows.meta,calc,babel,shapes.misc}
\RequirePackage[most]{tcolorbox}
\tcbuselibrary{skins,breakable}

\definecolor{coulDef}{HTML}{1F4E79}
\definecolor{coulProp}{HTML}{7030A0}
\definecolor{coulRem}{HTML}{C55A11}
\definecolor{coulEx}{HTML}{2E7D32}
\definecolor{coulExo}{HTML}{B03A2E}
\definecolor{coulPart}{HTML}{1F4E79}
\definecolor{coulMeth}{HTML}{1B6E4F}
\definecolor{coulCourbe}{HTML}{0B5ED7}
\definecolor{coulCourbeB}{HTML}{C00000}
\definecolor{coulCourbeC}{HTML}{2E7D32}
\definecolor{coulGrille}{HTML}{8C8C8C}
\definecolor{coulRep}{HTML}{1B6E4F}
\definecolor{coulBar}{HTML}{2E7D32}

\newcommand{\R}{\mathbb{R}}
\newcommand{\Rs}{\mathbb{R}^{*}}
\renewcommand{\le}{\leqslant}
\renewcommand{\ge}{\geqslant}
\newcommand{\vect}[1]{\overrightarrow{#1}}
\newcommand{\norme}[1]{\left\lVert #1 \right\rVert}
\newcommand{\intervalle}[2]{\left[#1\,;#2\right]}
\RequirePackage{cancel}
\renewcommand{\CancelColor}{\color{coulExo}}
\newcommand{\fois}[1]{\times{\color{coulExo}#1}}
\newcommand{\barre}[1]{\cancel{#1}}

\tikzset{grille/.style={coulGrille,line width=0.4pt}}
\newcommand{\quadrillage}[4]{\draw[grille,step=1] (#1,#2) grid (#3,#4);}
\tikzset{
  croix/.style={cross out,draw,line width=0.6pt,inner sep=0pt,minimum size=4pt},
  croixdroite/.style={croix,rotate=45,line width=0.8pt},
}
\newcommand{\pt}[3]{\path (#1) node[croix]{} node[#2,font=\small,inner sep=2.5pt]{$#3$};}

\newcommand{\lycee@repere}[4]{%
  \draw[grille,step=1] (#1,#2) grid (#3,#4);
  \draw[-{Stealth[length=2mm]},thick] (#1,0)--({#3+0.4},0) node[below left=-1pt]{$x$};
  \draw[-{Stealth[length=2mm]},thick] (0,#2)--(0,{#4+0.4}) node[below left=-1pt]{$y$};
  \node[below left,font=\scriptsize,inner sep=1.5pt] at (0,0) {$0$};}
\newcommand{\lycee@gradx}[1]{\foreach \k/\l in {#1}{%
  \draw (\k,0.07)--(\k,-0.07) node[below,font=\scriptsize,inner sep=1.5pt]{$\l$};}}
\newcommand{\lycee@grady}[1]{\foreach \k/\l in {#1}{%
  \draw (0.07,\k)--(-0.07,\k) node[left,font=\scriptsize,inner sep=1.5pt]{$\l$};}}
\newcommand{\lycee@intff}[2]{\left[#1\,;#2\right]}
\newcommand{\lycee@intfo}[2]{\left[#1\,;#2\right[}
\newcommand{\lycee@intof}[2]{\left]#1\,;#2\right]}
\newcommand{\lycee@intoo}[2]{\left]#1\,;#2\right[}
\newcommand{\lycee@ens}[1]{\left\{#1\right\}}
\AtBeginDocument{%
  \providecommand{\repere}{\lycee@repere}%
  \providecommand{\gradx}{\lycee@gradx}%
  \providecommand{\grady}{\lycee@grady}%
  \providecommand{\Cf}{\mathcal{C}\sb{\scriptscriptstyle f}}%
  \providecommand{\Cg}{\mathcal{C}\sb{\scriptscriptstyle g}}%
  \providecommand{\intff}{\lycee@intff}%
  \providecommand{\intfo}{\lycee@intfo}%
  \providecommand{\intof}{\lycee@intof}%
  \providecommand{\intoo}{\lycee@intoo}%
  \providecommand{\ens}{\lycee@ens}}

\newlist{questions}{enumerate}{3}
\setlist[questions,1]{label=\arabic*\textdegree),leftmargin=*,itemsep=2pt,topsep=3pt}
\setlist[questions,2]{label=\alph*),leftmargin=*,itemsep=1pt,topsep=2pt}
\setlist[questions,3]{label=\roman*),leftmargin=*,itemsep=1pt,topsep=1pt}
\setlist[itemize]{label=\textbullet,leftmargin=*,itemsep=2pt,topsep=3pt}

\newcommand{\besoinplace}[1]{\par\penalty\@M\begingroup
  \dimen@=#1\relax\dimen@ii=\pagegoal\advance\dimen@ii-\pagetotal
  \ifdim\dimen@>\dimen@ii\ifdim\dimen@ii>\z@\vfil\fi\break\fi\endgroup}

\newcommand{\lycee@pied}{\thepage/\pageref{LastPage}}
\newcommand{\entete}[2]{%
  \pagestyle{fancy}\fancyhf{}%
  \fancyhead[L]{\footnotesize\color{coulPart}\textbf{#1}}%
  \fancyhead[R]{\footnotesize\color{coulPart}\textbf{#2}}%
  \fancyfoot[C]{\footnotesize\lycee@pied}%
  \renewcommand{\headrulewidth}{0.4pt}%
  \renewcommand{\headrule}{\hbox to\headwidth{%
    \color{coulPart!40}\leaders\hrule height \headrulewidth\hfill}}}

\RequirePackage{ccicons}
\newcommand{\lycee@licence}{{\scriptsize\color{gray}%
  \copyright\ Luc Destombe --- \ccbyncsa\ CC BY-NC-SA 4.0}}
\newif\iflycee@licence \lycee@licencetrue
\newcommand{\sanslicence}{\lycee@licencefalse}
\AtBeginDocument{\iflycee@licence\AddToHookNext{shipout/foreground}{%
  \put(0,0){\rlap{\kern\dimexpr1in+\hoffset+\oddsidemargin+\textwidth\relax
    \llap{\raisebox{-\dimexpr1in+\voffset+\topmargin+\headheight+\headsep
      +\textheight+\footskip\relax}{\lycee@licence}}}}}\fi}

\newcommand{\formulesserrees}{%
  \setlength{\abovedisplayskip}{3pt}\setlength{\belowdisplayskip}{3pt}%
  \setlength{\abovedisplayshortskip}{2pt}\setlength{\belowdisplayshortskip}{3pt}}

\setlength{\parindent}{0pt}

\RequirePackage[implicit=false,hidelinks,pdfencoding=auto,psdextra,bookmarksopen,
  bookmarksopenlevel=1,pdfcreator={},pdfproducer={}]{hyperref}
\RequirePackage{bookmark}
\hypersetup{pdfauthor={Luc Destombe},pdfkeywords={CC BY-NC-SA 4.0}}
\newcounter{lycee@signet}
\newcommand{\signet}[3][]{\stepcounter{lycee@signet}%
  \edef\lycee@dest{\if\relax\detokenize{#1}\relax
    signet.\arabic{lycee@signet}\else#1\fi}%
  \hypertarget{\lycee@dest}{}\bookmark[level=#2,dest=\lycee@dest]{#3}}
\pdfstringdefDisableCommands{%
  \def\R{R}\def\vect#1{#1}\def\Cf{Cf}\def\Cg{Cg}\def\fois#1{×#1}%
  \def\barre#1{#1}\def\intervalle#1#2{[#1;#2]}\def\intff#1#2{[#1;#2]}%
  \def\textsuperscript#1{#1}\def\dfrac#1#2{#1/#2}\def\frac#1#2{#1/#2}%
  \def\vec#1{#1⃗}}%
\geometry{top=2cm,bottom=1cm,left=1cm,right=1cm,headheight=14pt,footskip=0.6cm}

\RequirePackage{pgfplots}
\pgfplotsset{compat=1.18}
\RequirePackage{tkz-tab}

\newcounter{partiecours}
\renewcommand{\thepartiecours}{\Roman{partiecours}}
\newcommand{\partiecours}[1]{%
  \refstepcounter{partiecours}%
  \Needspace*{8\baselineskip}%
  \bigskip
  \noindent\begin{tcolorbox}[enhanced,colback=coulPart,colframe=coulPart,
      boxrule=0pt,arc=2pt,left=8pt,right=8pt,top=4pt,bottom=4pt,
      before skip=6pt,after skip=10pt]
    \color{white}\bfseries\large \thepartiecours{} --- #1\signet{1}{\thepartiecours{} --- #1}
  \end{tcolorbox}%
  \addcontentsline{toc}{section}{\thepartiecours{} --- #1}%
}

\newtcolorbox{boitecours}[3][]{%
  enhanced,breakable,
  colback=white,colframe=#2,coltitle=white,
  fonttitle=\bfseries,
  attach boxed title to top left={xshift=8pt,yshift=-\tcboxedtitleheight/2},
  boxed title style={colback=#2,arc=2pt,boxrule=0pt},
  boxrule=0.9pt,arc=2pt,
  left=8pt,right=8pt,top=8pt,bottom=6pt,
  title={#3},#1}

\newcounter{methode}
\newenvironment{definitionbox}[1][Définition]%
  {\begin{boitecours}[phantom={\signet{2}{#1}}]{coulDef}{#1}}{\end{boitecours}}
\newenvironment{proprietebox}[1][Propriété]%
  {\begin{boitecours}[phantom={\signet{2}{#1}}]{coulProp}{#1}}{\end{boitecours}}
\newenvironment{methodebox}[1][Méthode]{\stepcounter{methode}%
  \begin{boitecours}[phantom={\signet[methode-\arabic{methode}]{2}{#1}}]{coulMeth}{#1}}%
  {\end{boitecours}}

\newcommand{\etape}[1]{\par\smallskip\textbf{\color{coulMeth}#1}\par\nobreak\smallskip}
\newcommand{\enonce}[1]{\emph{#1}\par\smallskip}
\newcommand{\reponse}[1]{\par\smallskip
  {\footnotesize\itshape\color{black!55}Réponses : #1}}
\newcommand{\demo}[1]{\textbf{\color{coulMeth}Démonstration.} #1}
\newcommand{\afaire}[1]{%
  \par\smallskip
  \noindent{\small\color{coulExo}$\blacktriangleright$\ \textbf{À faire :}
    fiche d'exercices du chapitre, #1.}\par\smallskip}

\newenvironment{calculs}{\setlength{\jot}{5pt}\@fleqntrue\@mathmargin\retraitcalculs
  \setlength{\abovedisplayskip}{4pt}\setlength{\belowdisplayskip}{4pt}%
  \setlength{\abovedisplayshortskip}{4pt}\setlength{\belowdisplayshortskip}{4pt}%
  \csname alignat*\endcsname{2}}%
  {\csname endalignat*\endcsname}
\newcommand{\rem}[1]{\text{\small\itshape\color{black!60}#1}}
\newcommand{\explique}[2]{\par\smallskip\noindent
  \begin{minipage}[t]{0.57\linewidth}\raggedright #1\end{minipage}\hfill
  \begin{minipage}[t]{0.40\linewidth}\raggedright\leavevmode
    {\small\itshape\color{black!60}#2}\end{minipage}\par}
\newcommand{\cas}[1]{\par\smallskip\noindent\textbf{\color{coulExo}#1}\nobreak}
\newcommand{\calc}[1]{\par\smallskip\textbf{\color{coulExo}#1}\par\nobreak}

\newtcolorbox{remarquebox}[1][Remarque]{%
  enhanced,breakable,colback=white,colframe=coulRem,
  boxrule=0pt,leftrule=2.5pt,arc=0pt,outer arc=0pt,
  left=8pt,right=6pt,top=5pt,bottom=5pt,
  before upper={\textbf{\color{coulRem}#1}\par\smallskip}}

\newtcolorbox{exemplebox}[1][Exemple]{%
  enhanced,breakable,colback=white,colframe=coulEx,
  boxrule=0pt,leftrule=2.5pt,arc=0pt,outer arc=0pt,
  left=8pt,right=6pt,top=5pt,bottom=5pt,
  before upper={\textbf{\color{coulEx}#1}\par\smallskip}}

\pgfplotsset{
  repere/.style={
    axis lines=middle,
    axis line style={-{Stealth[length=2.4mm]},thick},
    every axis x label/.style={at={(ticklabel* cs:1.0)},anchor=west},
    every axis y label/.style={at={(ticklabel* cs:1.0)},anchor=south},
    label style={font=\footnotesize},
    tick label style={font=\scriptsize},
    grid=both,
    major grid style={grille},
    minor tick num=0,
    tick align=inside,
    clip mode=individual,
  },
}
\tikzset{
  courbe/.style={coulCourbe,line width=1.1pt,smooth},
  courbeB/.style={coulCourbeB,line width=1.1pt,smooth},
  courbeC/.style={coulCourbeC,line width=1.1pt,smooth},
}

\newlist{etapes}{enumerate}{1}
\setlist[etapes]{label=\textbf{\arabic*.},leftmargin=*,itemsep=1pt,topsep=2pt}

\newlist{expressions}{enumerate}{1}
\setlist[expressions]{label={\Alph*\ $=$},leftmargin=*,itemsep=3pt,topsep=3pt}

\newcounter{exo}
\renewcommand{\theexo}{\ifnum\value{exo}<10 0\fi\arabic{exo}}
\newtcolorbox{exercice}[1][]{%
  enhanced,breakable,
  colback=white,colframe=coulExo,
  boxrule=0.9pt,arc=2pt,
  left=8pt,right=8pt,top=8pt,bottom=6pt,
  coltitle=white,fonttitle=\bfseries,
  attach boxed title to top left={xshift=8pt,yshift=-\tcboxedtitleheight/2},
  boxed title style={colback=coulExo,arc=2pt,boxrule=0pt},
  phantom={\signet{2}{Exercice \arabic{exo}}},
  title={Exercice \theexo\ifx\relax#1\relax\else{} --- #1\fi}}
\newenvironment{exo}[1][\relax]{\refstepcounter{exo}\begin{exercice}[#1]}{\end{exercice}}

  \RequirePackage{comment}
  \excludecomment{correction}
  \renewcommand{\reponse}[1]{}
  \newcommand{\autableau}{\hfill{\footnotesize\itshape\color{black!45}(au tableau)}}
  \newcommand{\etapeexemples}{%
    \par\smallskip\textbf{\color{coulMeth}Exemples}\autableau\par\nobreak\smallskip}
  \newcommand{\etapetableau}[1]{%
    \par\smallskip\textbf{\color{coulMeth}#1}\autableau\par\nobreak\smallskip}
  \newtcolorbox{exempletableau}[1][Exemple]{%
    enhanced,breakable,colback=white,colframe=coulEx,
    boxrule=0pt,leftrule=2.5pt,arc=0pt,outer arc=0pt,
    left=8pt,right=6pt,top=4pt,bottom=4pt,
    before skip=5pt,after skip=5pt,
    before upper={\textbf{\color{coulEx}#1}\autableau\par\smallskip}}

\newcommand{\coursEntete}{}
\newcommand{\coursNiveau}{}
\newcommand{\coursTitre}{}
\newcommand{\coursSurTitre}{}
\newcommand{\coursSousTitre}{}

\newcommand{\enteteCours}[1]{\renewcommand{\coursEntete}{#1}}
\newcommand{\chapitrecours}[2]{\enteteCours{Chapitre #1 --- #2}}
\newcommand{\niveaucours}[1]{\renewcommand{\coursNiveau}{#1}}
\newcommand{\titrecours}[3][]{%
  \renewcommand{\coursTitre}{#2}%
  \renewcommand{\coursSurTitre}{#1}%
  \renewcommand{\coursSousTitre}{#3}}

\pagestyle{fancy}
\fancyhf{}
\fancyhead[L]{\footnotesize\color{coulPart}\textbf{\coursEntete}}
\fancyhead[R]{\footnotesize\color{coulPart}\textbf{\coursNiveau}}
\fancyfoot[C]{\footnotesize\thepage}
\renewcommand{\headrulewidth}{0.4pt}
\renewcommand{\headrule}{\hbox to\headwidth{\color{coulPart!40}\leaders\hrule height \headrulewidth\hfill}}

\setlength{\parskip}{2pt}

\AtBeginDocument{%
  \begin{center}
    \begin{tcolorbox}[enhanced,width=0.72\linewidth,colback=white,colframe=coulPart,
        boxrule=1.6pt,arc=3pt,halign=center,top=10pt,bottom=10pt]
      \ifx\coursSurTitre\empty
        {\Huge\bfseries\color{coulPart} \coursTitre}\\[4pt]
      \else
        {\Huge\bfseries\color{coulPart} \coursTitre}\\[3pt]
        {\Large\bfseries\color{coulPart} \coursSurTitre}\\[5pt]
      \fi
      {\normalsize \coursSousTitre}%
      \\[2pt]{\small\itshape Version élève}
    \end{tcolorbox}
  \end{center}%
}
\makeatother

\chapitrecours{1}{Fractions et puissances}
\niveaucours{Seconde}
\titrecours{FRACTIONS ET PUISSANCES}{Cours --- Seconde}

\begin{document}

\begin{remarquebox}[Comment utiliser ce cours]
Chaque partie commence par ce qu'il faut \textbf{savoir} (définitions, propriétés), puis
vient ce qu'il faut \textbf{savoir faire} : une méthode, avec des
\textbf{exemples} faits en classe, puis des exercices
\og\textbf{À vous de jouer}\fg{} à chercher seul.

\smallskip
Sur cette feuille, seuls les \textbf{énoncés} des exemples figurent : la résolution,
signalée par la mention \emph{(au tableau)}, est à rédiger dans le cahier.

\end{remarquebox}

\partiecours{Nombres relatifs et priorités}

\begin{definitionbox}[Définition 1 --- Les sens du signe moins]
Le signe $-$ a trois sens :
\begin{itemize}
  \item une \textbf{soustraction} : $9-2$ se lit \og neuf moins deux\fg{} ;
  \item un \textbf{nombre négatif} : $-6$ se lit \og moins six\fg{} ;
  \item l'\textbf{opposé} : $-(-4)$ se lit \og l'opposé de moins quatre\fg{}, et vaut $4$.
\end{itemize}
\end{definitionbox}

\begin{proprietebox}[Propriété 1 --- Règles des signes]
\begin{minipage}[t]{0.48\linewidth}
\textbf{Somme}
\begin{itemize}
  \item deux nombres de même signe : on ajoute leurs distances à zéro et on garde le
        signe : $-5-3=-8$ ;
  \item de signes contraires : on soustrait les distances à zéro, le signe est celui
        du plus éloigné de zéro : $-5+3=-2$.
\end{itemize}
\end{minipage}\hfill
\begin{minipage}[t]{0.48\linewidth}
\textbf{Produit et quotient}
\begin{itemize}
  \item deux nombres de même signe : le résultat est positif :
        $(-5)\times(-3)=15$ ;
  \item de signes contraires : le résultat est négatif : $(-6)\div 3=-2$.
\end{itemize}
\end{minipage}

\smallskip
Soustraire un nombre, c'est ajouter son opposé : $4-(-7)=4+7$.
\end{proprietebox}

\begin{methodebox}[Méthode 1 --- Calculer avec des nombres relatifs]
\etapeexemples
Calculer :
\begin{questions}[label=\alph*)]
  \item $A=-7-9$
  \item $B=15-(-6)$
  \item $C=3\times(-8)$
  \item $D=\dfrac{-(-18)}{-6}$
\end{questions}

\etape{À vous de jouer}
Calculer :
\[
  E=-5-12 \qquad F=8-(-14) \qquad G=-4\times(-9) \qquad H=\dfrac{-(-35)}{-7}
  \qquad I=-6-(-10)
\]

\end{methodebox}

\begin{proprietebox}[Propriété 2 --- Priorités de calcul]
On effectue, dans cet ordre :
\begin{enumerate}[label=\arabic*.]
  \item les calculs entre \textbf{parenthèses}, en commençant par les plus intérieures ;
  \item les \textbf{puissances} ;
  \item les \textbf{multiplications et divisions}, de gauche à droite ;
  \item les \textbf{additions et soustractions}, de gauche à droite.
\end{enumerate}
\end{proprietebox}

\begin{methodebox}[Méthode 2 --- Respecter les priorités de calcul]
\etapeexemples
Calculer, en écrivant une étape par ligne :
\begin{questions}[label=\alph*)]
  \item $A=6^{2}-5\times(-3)$
  \item $B=-8-2^{4}$
  \item $C=7+(-10+4)\times 3$
  \item $D=-24\div(-4)\times 3^{2}$
\end{questions}

\begin{remarquebox}[Erreur à éviter]
Dans $D$, on ne commence pas par $(-4)\times 9$ : la division et la multiplication ont
la même priorité, on les fait dans l'ordre où elles sont écrites.
\end{remarquebox}

\etape{À vous de jouer}
Calculer :
\[
  E=4^{2}-2\times(-7) \qquad F=-12-3^{2} \qquad G=9+(-6+1)\times 2 \qquad
  H=-30\div(-5)\times 2^{3}
\]

\end{methodebox}

\afaire{exercices 1 à 3}

\partiecours{Additionner et soustraire des fractions}

\begin{proprietebox}[Propriété 3 --- Le signe d'une fraction]
Pour $b\ne 0$ : \quad $\dfrac{-a}{b}=\dfrac{a}{-b}=-\dfrac{a}{b}$ \quad et \quad
$\dfrac{-a}{-b}=\dfrac{a}{b}$.

\smallskip
On écrit le signe moins d'une fraction au numérateur ou devant le trait de fraction,
jamais au dénominateur. Un signe moins devant une fraction s'applique à \textbf{tout}
le numérateur : on le met entre parenthèses.
\end{proprietebox}

\begin{exemplebox}[Exemple]
\begin{calculs}
3-\dfrac{x-1}{5} &= \dfrac{3\fois{5}}{1\fois{5}}-\dfrac{x-1}{5} &\qquad& \rem{$3=\frac31$, écrit avec le dénominateur $5$}\\
3-\dfrac{x-1}{5} &= \dfrac{15}{5}-\dfrac{x-1}{5} &&\\
3-\dfrac{x-1}{5} &= \dfrac{15-(x-1)}{5} &\qquad& \rem{parenthèses : le moins porte sur $x$ et sur $-1$}\\
3-\dfrac{x-1}{5} &= \dfrac{16-x}{5} &\qquad& \rem{$-(x-1)=-x+1$}
\end{calculs}
\end{exemplebox}

\begin{proprietebox}[Propriété 4 --- Somme de fractions]
Pour additionner (ou soustraire) des fractions, on les écrit avec le \textbf{même
dénominateur}, puis on additionne (ou on soustrait) les numérateurs ; le dénominateur
ne change pas :
\[
  \dfrac{a}{d}+\dfrac{b}{d}=\dfrac{a+b}{d} \qquad\qquad \dfrac{a}{d}-\dfrac{b}{d}=\dfrac{a-b}{d}
\]
\end{proprietebox}

\begin{methodebox}[Méthode 3 --- Additionner et soustraire des fractions]
\etapeexemples
Calculer et donner le résultat sous forme simplifiée :
\begin{questions}[label=\alph*)]
  \item $A=\dfrac{-7}{-6}+\dfrac{-3}{12}$
  \item $B=\dfrac{4}{5}-\dfrac{-2}{3}$
  \item $C=\dfrac{1}{3}-\dfrac{1}{6}+\dfrac{1}{9}-\dfrac{1}{18}$
\end{questions}

\etape{À vous de jouer}
Calculer et donner le résultat sous forme simplifiée :
\[
  D=\dfrac{-2}{-3}+\dfrac{-5}{6} \qquad
  E=\dfrac{5}{6}-\dfrac{-3}{4} \qquad
  F=\dfrac{1}{2}-\dfrac{1}{4}+\dfrac{1}{10} \qquad
  G=4-\dfrac{7}{3}
\]

\end{methodebox}

\afaire{exercices 4 à 7}

\partiecours{Multiplier et diviser des fractions}

\begin{proprietebox}[Propriété 5 --- Produit et quotient]
Pour $b$, $c$ et $d$ non nuls :
\[
  \dfrac{a}{b}\times\dfrac{c}{d}=\dfrac{a\times c}{b\times d}
  \qquad\qquad
  \dfrac{a}{b}\div\dfrac{c}{d}=\dfrac{a}{b}\times\dfrac{d}{c}
\]
Diviser par une fraction, c'est multiplier par son \textbf{inverse}.
\end{proprietebox}

\begin{methodebox}[Méthode 4 --- Multiplier et diviser des fractions]
\etapeexemples
Calculer et donner le résultat sous forme simplifiée :
\begin{questions}[label=\alph*)]
  \item $A=\dfrac{3}{-4}\times\dfrac{-5}{7}$
  \item $B=6\times\dfrac{4}{5}$
  \item $C=\dfrac{5}{6}\div\dfrac{-10}{9}$
\end{questions}

\etape{À vous de jouer}
Calculer et donner le résultat sous forme simplifiée :
\[
  D=\dfrac{2}{-7}\times\dfrac{-3}{5} \qquad
  E=4\times\dfrac{2}{9} \qquad
  F=\dfrac{3}{10}\div\dfrac{-9}{5} \qquad
  G=\dfrac{-4}{9}\div 8
\]

\end{methodebox}

\begin{remarquebox}[Simplifier : seulement des produits]
Dans un quotient de \textbf{produits}, on peut barrer un même facteur en haut et en
bas. Jamais dans une \textbf{somme} : dans $\frac{-2+9}{5+9}$, on ne barre pas les $9$.
Quand on barre tout le numérateur, il reste $1$, et non $0$.
\end{remarquebox}

\begin{methodebox}[Méthode 5 --- Simplifier avant de calculer]
\etapeexemples
Simplifier, puis calculer :
\begin{questions}[label=\alph*)]
  \item $A=\dfrac{-9\times 10}{6\times 25}$
  \item $B=\dfrac{3\times(-4)\times(-5)}{(-6)\times 7}$
  \item $C=\dfrac{-2+9}{5+9}$
  \item $D=\dfrac{28}{15}\times\dfrac{25}{14}$
  \item $E=\dfrac{6\times 5}{10\times 9}$
\end{questions}

\etape{À vous de jouer}
Simplifier, puis calculer :
\[
  F=\dfrac{-8\times 15}{12\times 10} \qquad
  G=\dfrac{6+4}{5+4} \qquad
  H=\dfrac{5\times(-2)\times 3}{3\times 5} \qquad
  I=\dfrac{21}{20}\times\dfrac{8}{35} \qquad
  J=\dfrac{18}{35}\times\dfrac{14}{27}\times\dfrac{5}{4}
\]

\end{methodebox}

\afaire{exercices 8 à 11}

\partiecours{Fractions et priorités de calcul}

\begin{proprietebox}[Propriété 6 --- Fraction de fractions]
Pour $b$, $c$ et $d$ non nuls : \quad
$\dfrac{\;\dfrac{a}{b}\;}{\;\dfrac{c}{d}\;}=\dfrac{a}{b}\div\dfrac{c}{d}=\dfrac{a}{b}\times\dfrac{d}{c}$.

\smallskip
Quand le numérateur ou le dénominateur est un calcul, on le calcule à part, puis on
divise.
\end{proprietebox}

\begin{methodebox}[Méthode 6 --- Mener un calcul avec des fractions]
\etapeexemples
Calculer et donner le résultat sous forme simplifiée :
\begin{questions}[label=\alph*)]
  \item $A=\dfrac{5}{6}-\dfrac{-2}{3}\times\dfrac{3}{2^{2}}$
  \item $B=\dfrac{-5}{\;3+\dfrac{3}{2^{2}}\;}$
  \item $C=\left(\dfrac{1}{2}+\dfrac{1}{3}\right)\div\left(\dfrac{3}{4}-\dfrac{1}{6}\right)$
  \item $D=\dfrac{1}{2}+\dfrac{1}{3}\div\dfrac{3}{4}-\dfrac{1}{6}$
\end{questions}

\etape{À vous de jouer}
Calculer et donner le résultat sous forme simplifiée :
\[
  E=\left(2+\dfrac{1}{3}-\dfrac{3}{4}\right)\times\left(\dfrac{1}{2}-\dfrac{2}{3}\right)
  \qquad
  F=\dfrac{\;\dfrac{1}{3}+\dfrac{1}{4}\;}{\;\dfrac{5}{6}-\dfrac{1}{4}+\dfrac{1}{3}\;}
\]
\[
  G=-\dfrac{3}{7}\times\left[\left(\dfrac{1}{2}+\dfrac{5}{4}\right)-\left(\dfrac{1}{3}-\dfrac{1}{6}\right)\right]
  \qquad
  H=\left(\dfrac{-3}{4}+\dfrac{1}{2}\right)\div\left(\dfrac{2}{3}+\dfrac{5}{-6}\right)
\]

\end{methodebox}

\afaire{exercices 12 à 14}

\partiecours{Produit en croix}

\begin{proprietebox}[Propriété 7 --- Produit en croix]
Pour $b\ne 0$ et $d\ne 0$ : \quad si $\dfrac{a}{b}=\dfrac{c}{d}$, alors
$a\times d=b\times c$.
\end{proprietebox}

\begin{remarquebox}[Pourquoi ?]
On écrit les deux fractions avec le dénominateur commun $b\times d$ :
$\dfrac{a\times d}{b\times d}=\dfrac{c\times b}{d\times b}$. Deux fractions égales de même
dénominateur ont le même numérateur : $a\times d=b\times c$.
\end{remarquebox}

\begin{methodebox}[Méthode 7 --- Résoudre une équation par le produit en croix]
\etapeexemples
Résoudre les équations suivantes :
\begin{questions}[label=\alph*)]
  \item $\dfrac{x}{3}=\dfrac{5}{4}$
  \item $\dfrac{x}{4}=\dfrac{x+3}{6}$
  \item $\dfrac{x+3}{2}=\dfrac{2x-1}{3}$
  \item $\dfrac{2x+1}{4}=\dfrac{3x-2}{-2}$
\end{questions}

\begin{remarquebox}[Vérifier une solution]
On remplace $x$ par la valeur trouvée dans chaque membre. Pour b) :
$\frac{6}{4}=\frac32$ et $\frac{6+3}{6}=\frac96=\frac32$ ; les deux membres sont égaux,
donc $6$ est bien la solution.
\end{remarquebox}

\etape{À vous de jouer}
Résoudre, puis vérifier :
\[
  \text{e) } \dfrac{3}{7}=\dfrac{x}{14} \qquad
  \text{f) } \dfrac{3x}{4}=\dfrac{7-x}{2} \qquad
  \text{g) } \dfrac{2x+1}{3}=\dfrac{x+4}{2} \qquad
  \text{h) } \dfrac{3x-2}{4}=\dfrac{5-x}{8} \qquad
  \text{i) } \dfrac{5x+2}{3}=\dfrac{4-x}{-2}
\]

\end{methodebox}

\afaire{exercices 15 à 17}

\begin{methodebox}[Méthode 8 --- Résoudre une équation à plusieurs fractions]
Quand un membre contient une somme de fractions, le produit en croix ne s'applique
pas : on multiplie \textbf{chaque terme} des deux membres par un dénominateur commun.
\etapeexemples
Résoudre les équations suivantes :
\begin{questions}[label=\alph*)]
  \item $x+\dfrac{x}{3}=12$
  \item $\dfrac{t}{4}-\dfrac{t+1}{6}=1$
  \item $\dfrac{z-1}{2}-\dfrac{z-3}{6}=\dfrac{2z+1}{9}$
\end{questions}

\etape{À vous de jouer}
Résoudre :
\[
  \text{d) } \dfrac{x}{3}-\dfrac{x}{5}=4 \qquad
  \text{e) } \dfrac{3t-1}{2}+\dfrac{t-3}{4}=t \qquad
  \text{f) } \dfrac{y}{2}+\dfrac{2y}{3}=\dfrac{y}{4}+11
\]
\[
  \text{g) } \dfrac{y+6}{4}=\dfrac{y-2}{3}-\dfrac{y}{6} \qquad
  \text{h) } \dfrac{3z+1}{4}-\dfrac{5z-2}{3}=\dfrac{1-13z}{12}
\]

\end{methodebox}

\afaire{exercices 18 et 19}

\partiecours{Fractions et problèmes}

\begin{proprietebox}[Propriété 8 --- Fraction d'une quantité]
Prendre les $\dfrac{a}{b}$ \textbf{de} quelque chose, c'est \textbf{multiplier} par
$\dfrac{a}{b}$ : les $\dfrac23$ de $60$ font $\dfrac23\times 60=40$.

\smallskip
Une fraction prise \og du reste\fg{} porte sur ce qui reste, pas sur le total : on
calcule d'abord ce reste.
\end{proprietebox}

\begin{methodebox}[Méthode 9 --- Résoudre un problème de fractions]
\etapeexemples
\begin{questions}[label=\alph*)]
  \item Un train de $360$ places part de Lille avec les trois quarts de ses places
        occupées. À Arras, des voyageurs occupent les deux tiers des places restées
        libres. Combien de places sont encore libres ?
  \item Malo a un forfait mobile de $60$~Go. Il en utilise le quart pour les vidéos,
        puis les deux cinquièmes de ce qui reste pour les réseaux sociaux. Quelle
        fraction du forfait lui reste-t-il ? Combien de gigaoctets cela fait-il ?
  \item Une boutique en ligne vend les deux tiers de son stock de casques audio le
        premier jour des soldes, puis les trois cinquièmes du reste le deuxième jour. Il
        reste alors $20$ casques. Combien en a-t-elle vendu ?
\end{questions}

\etape{À vous de jouer}
\begin{questions}
  \item Une cuve de $240$~L est remplie aux cinq sixièmes. On en vide les trois quarts.
        Combien de litres reste-t-il dans la cuve ?
  \item Sami dépense le quart de ses économies pour un téléphone, puis les deux tiers de
        ce qui reste pour un vélo. Il lui reste $75$~€. Combien avait-il économisé ?
\end{questions}

\end{methodebox}

\afaire{exercices 20 à 23}

\partiecours{Puissances}

\begin{definitionbox}[Définition 2 --- Puissance d'un nombre]
Pour un nombre $a$ et un entier $n\ge 1$ :
\[
  a^{n}=\underbrace{a\times a\times\dots\times a}_{n\text{ facteurs}}
  \qquad a^{0}=1 \qquad
  a^{-n}=\dfrac{1}{a^{n}} \quad (a\ne 0)
\]
Un exposant positif compte des multiplications, un exposant négatif des divisions :
$5^{-2}=\dfrac{1}{5\times 5}=\dfrac{1}{25}$.

\smallskip
En particulier, $a^{-1}=\dfrac1a$ : $a^{-1}$ est l'\textbf{inverse} de $a$.
\end{definitionbox}

\begin{definitionbox}[Définition 3 --- Puissances de 10]
$10^{n}=1\underbrace{0\dots0}_{n\text{ zéros}}$ \qquad et \qquad
$10^{-n}=\dfrac{1}{10^{n}}=0{,}\underbrace{0\dots0}_{n-1\text{ zéros}}1$
\qquad ($n$ chiffres après la virgule).

\smallskip
Exemples : $10^{6}=1~000~000$ \quad ; \quad $10^{-4}=0{,}0001$.
\end{definitionbox}

\begin{remarquebox}[Signe d'une puissance]
Un nombre négatif élevé à une puissance \textbf{paire} donne un résultat positif ; à
une puissance \textbf{impaire}, un résultat négatif : $(-2)^{4}=16$, $(-2)^{3}=-8$.
Attention aux parenthèses : $(-3)^{2}=9$, mais $-3^{2}=-(3\times 3)=-9$.
\end{remarquebox}

\begin{methodebox}[Méthode 10 --- Calculer avec des puissances]
\etapeexemples
Calculer :
\begin{questions}[label=\alph*)]
  \item $A=(3^{2}-4)\times(-1)^{7}$
  \item $B=(-2)^{3}+(-3)^{2}\times(-2)$
  \item $C=-3^{2}+(-3)^{2}-(-3)^{3}$
\end{questions}

\etape{À vous de jouer}
Calculer :
\[
  D=(2^{3}-5)\times(-1)^{6} \qquad E=(-1)^{9}+(-2)^{4}\times(-3) \qquad
  F=(-6)^{2}-6^{2}+(-2)^{3}
\]

\end{methodebox}

\begin{methodebox}[Méthode 11 --- Passer d'un exposant négatif à une fraction]
\etapeexemples
\begin{questions}[label=\alph*)]
  \item Écrire sous forme de fraction : $A=(-3)^{-2}$ et $B=-4^{-2}$.
  \item Écrire sous la forme $a^{n}$ : $C=\dfrac{1}{9}$, \;
        $D=\dfrac{1}{2\times 2\times 2\times 2\times 2\times 2\times 2}$ et
        $E=\dfrac{1}{(-4)\times(-4)\times(-4)}$.
\end{questions}

\etape{À vous de jouer}
\begin{questions}
  \item Écrire sous forme de fraction : \quad $F=(-2)^{-3}$ \qquad $G=-3^{-2}$ \qquad
        $H=\left(\dfrac{1}{5}\right)^{-1}$
  \item Écrire sous la forme $a^{n}$ : \quad $I=\dfrac{1}{6}$ \qquad
        $J=\dfrac{1}{5\times 5\times 5\times 5}$ \qquad $K=\dfrac{1}{(-7)\times(-7)}$
\end{questions}

\end{methodebox}

\afaire{exercices 24 à 29}

\partiecours{Règles de calcul sur les puissances}

\begin{proprietebox}[Propriété 9 --- Règles de calcul]
Pour $a$ et $b$ non nuls, $n$ et $p$ entiers relatifs :
\[
\renewcommand{\arraystretch}{2.2}
\begin{array}{lll}
  a^{n}\times a^{p}=a^{n+p} & \quad & \text{exemple : } 5^{4}\times 5^{-6}=5^{4+(-6)}=5^{-2}\\
  \dfrac{a^{n}}{a^{p}}=a^{n-p} & & \text{exemple : } \dfrac{(-2)^{9}}{(-2)^{4}}=(-2)^{9-4}=(-2)^{5}\\
  \left(a^{n}\right)^{p}=a^{n\times p} & & \text{exemple : } \left(4^{3}\right)^{5}=4^{3\times 5}=4^{15}\\
  (a\times b)^{n}=a^{n}\times b^{n} & & \text{exemple : } (3x)^{2}=3^{2}\times x^{2}=9x^{2}\\
  \left(\dfrac{a}{b}\right)^{n}=\dfrac{a^{n}}{b^{n}} & & \text{exemple : } \left(\dfrac{3}{5}\right)^{2}=\dfrac{3^{2}}{5^{2}}=\dfrac{9}{25}
\end{array}
\]
\end{proprietebox}

\begin{methodebox}[Méthode 12 --- Simplifier avec les règles sur les puissances]
\etapeexemples
Simplifier, puis calculer :
\begin{questions}[label=\alph*)]
  \item $A=\dfrac{(2^{2})^{-3}\times 3}{(2\times 3)^{-2}}$
  \item $B=\dfrac{3^{4}\times 10^{-5}\times(3^{2})^{3}\times(10^{-1})^{2}}{3^{8}\times 10^{-4}}$
\end{questions}

\etape{À vous de jouer}
Simplifier, puis calculer :
\[
  C=\dfrac{(3^{2})^{-2}\times 2}{(3\times 2)^{-3}} \qquad
  D=\dfrac{2^{5}\times 10^{-4}}{2^{3}\times 10^{-6}} \qquad
  E=\dfrac{a^{4}\times(a^{-3})^{2}}{(a^{2})^{3}} \quad (a\ne 0)
\]

\end{methodebox}

\afaire{exercices 30 à 34}

\partiecours{Notation scientifique}

\begin{definitionbox}[Définition 4 --- Notation scientifique]
La \textbf{notation scientifique} d'un nombre décimal non nul est son écriture
$a\times 10^{n}$, où $a$ a \textbf{un seul chiffre non nul avant la virgule} et $n$
est un entier relatif.

\smallskip
Exemples : $27~800=2{,}78\times 10^{4}$ \quad ; \quad
$-40~500~000=-4{,}05\times 10^{7}$ \quad ; \quad $0{,}000~000~63=6{,}3\times 10^{-7}$.
\end{definitionbox}

\begin{remarquebox}[Pour ne pas se tromper d'exposant]
Le produit $a\times 10^{n}$ ne doit pas changer de valeur : quand la virgule de $a$
recule de $k$ rangs (le nombre devient plus petit), l'exposant augmente de $k$ ; quand
elle avance, l'exposant diminue.
\end{remarquebox}

\begin{methodebox}[Méthode 13 --- Écrire un nombre en notation scientifique]
\etapeexemples
Donner la notation scientifique de :
\begin{questions}[label=\alph*)]
  \item $A=5~600~000$ \quad et \quad $B=0{,}000~47$
  \item $C=318{,}2\times 10^{-6}$
  \item $D=\dfrac{4\times 10^{2}\times 9\times 10^{5}}{60\times 10^{3}}$
  \item $E=6\times 10^{3}+47\times 10^{-2}$
\end{questions}

\etape{À vous de jouer}
Donner la notation scientifique de :
\[
  F=3~470~000 \qquad G=0{,}000~902 \qquad H=72{,}6\times 10^{-5} \qquad
  I=\dfrac{8\times 10^{5}\times 3\times 10^{-3}}{6\times 10^{4}}
\]

\end{methodebox}

\afaire{exercices 35 à 39}

\end{document}
